\section{Ejercicio 1}

\subsection{Consideraciones de diseño}
Se tomaron como reservados los exponentes 4'h0 y 4'hf para representaciones especiales:
\begin{itemize}
    \item Overflow: se representa $\pm\infty$ con todos los bits del exponente en 1 y la mantisa en 0.
    \item Underflow: se representa $\pm 0$ con todos los bits del exopnente en 0 y la mantisa en 0.
    \item NaN: si uno de los datos de entrada este se propaga a la salida. Si por algún motivo algún otro bit además del MSB de la mantisa está en 1, el mismo se filtra.
\end{itemize}
Además, se decidió no soportar el bit escondido en 0 (implica que el exponente esté todo en 0). El bit escondido se toma en 1 siempre.

\subsection{Arquitectura}


\subsection{Código}
\inputminted[fontsize=\footnotesize]{systemverilog}{../floating_point_multiplier/rtl/floating_point_multiplier.sv}

\subsection{Testbench}
\inputminted[fontsize=\footnotesize]{systemverilog}{../floating_point_multiplier/tb/tb_floating_point_multiplier.sv}
\subsubsection{Datasets}
\inputminted[fontsize=\footnotesize]{systemverilog}{../floating_point_multiplier/tb/data_in_range.sv}
\inputminted[fontsize=\footnotesize]{systemverilog}{../floating_point_multiplier/tb/data_overflow.sv}
\inputminted[fontsize=\footnotesize]{systemverilog}{../floating_point_multiplier/tb/data_underflow.sv}
